% Computational Experiments with Stein Factorization in Macaulay2
% Companion note to the SteinFactorizationM2 prototype.
% Build:  pdflatex IMPLEMENTATION.tex   (twice, for the cross-references)
\documentclass[11pt,reqno]{amsart}

\usepackage[T1]{fontenc}
\usepackage[utf8]{inputenc}
\usepackage{amsmath,amssymb,amsthm}
\usepackage{booktabs,longtable,array}
\usepackage{microtype}
\usepackage[margin=1in]{geometry}
\usepackage{xcolor}
\usepackage[colorlinks=true,linkcolor=blue!50!black,urlcolor=blue!50!black,
            citecolor=blue!50!black]{hyperref}

% Verbatim blocks hold Macaulay2 input and output.
\usepackage{fancyvrb}
\fvset{fontsize=\small,xleftmargin=1em}

% Long \texttt identifiers cannot hyphenate; give the paragraph builder room.
\setlength{\emergencystretch}{3em}

% Notation.
\newcommand{\PP}{\ensuremath{\mathbb{P}}}
\newcommand{\QQ}{\ensuremath{\mathbb{Q}}}
\newcommand{\RR}{\ensuremath{\mathbb{R}}}
\newcommand{\ZZ}{\ensuremath{\mathbb{Z}}}
\newcommand{\sO}{\ensuremath{\mathcal{O}}}

\DeclareMathOperator{\Hom}{Hom}
\DeclareMathOperator{\Spec}{Spec}
\DeclareMathOperator{\Proj}{Proj}
\DeclareMathOperator{\Bl}{Bl}

\theoremstyle{plain}
\newtheorem{theorem}{Theorem}[section]
\newtheorem{lemma}[theorem]{Lemma}
\newtheorem{proposition}[theorem]{Proposition}
\theoremstyle{definition}
\newtheorem{definition}[theorem]{Definition}
\newtheorem{remark}[theorem]{Remark}

\title{Computational Experiments with Stein Factorization in Macaulay2}
\author{Takehiko Yasuda}
\date{July 2026}

\begin{document}

\begin{abstract}
We implement a prototype of a Stein factorization algorithm in Macaulay2
and test it on several explicit projective maps. The implementation
follows the bigraded Hom-module construction in the author's algorithm:
\[C=\Hom_R(R_{\geq\mathbf{r}},R)_{(0,\geq 0)}\]
The module is used to reconstruct a graded coordinate algebra for the
Stein intermediate scheme and, where possible, the graph of the
connected-fiber morphism. This note is not intended as a formal
verification of the algorithm; rather, it records a reproducible
computational experiment that demonstrates the practical feasibility of
the construction on concrete examples. All calculations are included in
the repository and can be reproduced with Macaulay2 1.24.11 or later.
The code and this note were written essentially by an AI system, with
minor edits by the author; see \hyperref[sec:ai]{Use of AI}.
\end{abstract}

\maketitle


\section{Purpose and background}\label{sec:background}

For a projective morphism of algebraic varieties
\[f:Y\longrightarrow X,\]
a \emph{Stein factorization} is a decomposition
\[Y\xrightarrow{h} Z\xrightarrow{g} X\]
in which \(h\) is proper with \(h_*\sO_Y=\sO_Z\) (so that \(h\) has
connected fibers) and \(g\) is finite; concretely
\(Z=\Spec_X f_*\sO_Y\). The purpose of this project is to turn an
algorithm described by Yasuda (2026) into an executable Macaulay2
prototype and to see how it behaves on concrete examples from projective
geometry.

The input is the graph of \(f\) represented as a bigraded projective
scheme. If
\[R=S/I_{\Gamma_f},\]
then the variables in the source block have degrees \((\ast,0)\) and
those in the target block \((0,\ast)\), the starred entries being
positive; in the standard case they are \((1,0)\) and \((0,1)\), and
allowing larger values is what admits weighted projective spaces. What
matters is the split itself: it is what makes it possible to separate the
source direction from the target direction during the Hom-module
computation.

\section{The computational construction}\label{sec:construction}

For a sufficiently large bidegree \(\mathbf{r}=(r_1,r_2)\), the implementation
computes
\[C=\Hom_R(R_{\geq\mathbf{r}},R)_{(0,\geq0)}.\]
The truncation bound is obtained from shifts in a bounded bigraded free
resolution of \(R\) over the ambient polynomial ring. This is the
computational counterpart of the bound in Corollary 4.3 of
\hyperref[sec:references]{Yasuda (2026)}. The bigraded construction extends
the monograded global extension module computation of
\hyperref[sec:references]{Smith (2000)} to the bigraded setting.

For \(\mathbf{r}\) large enough, \(\Hom_R(R_{\geq\mathbf{r}},R)\) computes the
sections of the structure sheaf of \(\Gamma_f\cong Y\), so its
\((0,n)\)-part is
\[H^0(Y,f^*\sO_X(n))=H^0(X,f_*\sO_Y\otimes\sO_X(n))=H^0(Z,g^*\sO_X(n)).\]
These parts are the summands of \(C\): the subscript in the display above
is shorthand for the direct sum
\[C=\Hom_R(R_{\geq\mathbf{r}},R)_{(0,\geq0)}
  =\bigoplus_{n\geq0}\Hom_R(R_{\geq\mathbf{r}},R)_{(0,n)},\]
in which the first degree is held at \(0\) while the second runs over all
\(n\geq0\). A direct sum of graded pieces selected this way --- one
coordinate of the bigrading fixed, the other free --- is called a
\emph{strand}, and we call \(C\) the \((0,\geq0)\)-strand from here on.
(The same word is used in a different sense for the linear part of a free
resolution, which is not what is meant here.) So the
\((0,\geq0)\)-strand is the section ring of the Stein intermediate \(Z\)
with respect to \(g^*\sO_X(1)\). That polarization is pulled back from
\(X\), so it need not be \(\sO_Z(1)\): in every example below the finite
map \(g\) is a coordinatewise power map, and \(g^*\sO_X(1)=\sO_Z(d)\) for
some \(d>1\). Write
\[T=\bigoplus_{n\geq0}H^0(Z,\sO_Z(n))\]
for the section ring of \(\sO_Z(1)\) itself, so that
\(T=k[x_0,\dots,x_m]\) when \(Z=\PP^m\). What the strand computes is then
not \(T\) but
\[C=\bigoplus_{n\geq0}H^0(Z,\sO_Z(dn))=\bigoplus_{n\geq0}T_{dn}=T^{(d)},\]
the subring of \(T\) spanned by the graded pieces whose degree is a
multiple of \(d\) --- for \(Z=\PP^m\), the span of the monomials of degree
\(d,2d,3d,\dots\). This \(T^{(d)}\) is classically called the \(d\)-th
\emph{Veronese subring} of \(T\), since it is the coordinate ring of
\(Z\) in its \(d\)-uple Veronese embedding; the sections below refer to it
by that name.

There are three main functions, one for each step of the construction:
compute the truncated Hom module, read a coordinate algebra for \(Z\) off
it, and recover the graph of \(h\). Each returns a table that the next one
consumes, so they are meant to be called in this order.

\begin{verbatim}
homData     = steinHomData(S, Igraph)
algebraData = steinCoordinateAlgebra homData
graphData   = directSteinGraph(homData, algebraData)
\end{verbatim}

\subsubsection*{The Hom module ---
  \texorpdfstring{{\normalfont\ttfamily steinHomData}}{steinHomData}}

\texttt{S} is the ambient bigraded polynomial ring, the coordinate ring
of the product of the two (weighted) projective spaces in which the graph
sits; its variables are the two blocks described above.
\texttt{Igraph} is the bihomogeneous ideal \(I_{\Gamma_f}\subset S\)
cutting the graph out of that product, so that \(R=S/I_{\Gamma_f}\). 

% Those two arguments are all that is needed. Corollary 4.3 is stated in
% terms of the dimensions \(d_s\) of the two ambient (weighted) projective
% spaces and the sums \(c_s=\sum_{t=0}^{d_s}c_{s,t}\) of the degrees of the
% variables in each block, but the degrees of \texttt{S} already determine
% all four, so they are computed rather than passed. (\texttt{blockDegreeData
% S} returns them, and Section \ref{sec:weighted} uses it to exhibit the one
% case where they are not obvious: \(c_s\) is a weight sum, equal to the
% number of variables in the block only when every weight is \(1\).)

The returned \texttt{homData} carries everything the later steps need: the
bound \(\mathbf{r}\) and whether it is certified, the ring \(R\), the
truncation \(R_{\geq\mathbf{r}}\), the module
\(\Hom_R(R_{\geq\mathbf{r}},R)_{\geq(0,0)}\) together with the matrix of its
generators, and the indices of those generators lying in the
\((0,\geq0)\)-strand.

\subsubsection*{The coordinate algebra ---
  \texorpdfstring{{\normalfont\ttfamily steinCoordinateAlgebra}}{steinCoordinateAlgebra}}

The strand \(C\) consists of homomorphisms rather than of elements of
\(R\), so it carries no evident multiplication. It acquires one by being
realized inside a localization of \(R\): fix a nonzero
\(\gamma\in R_{\geq\mathbf{r}}\) and send
\[\varphi\longmapsto\varphi(\gamma)/\gamma\in R[1/\gamma].\]
Here \(R\) is a domain, \(\Gamma_f\) being a variety, so this is injective
and does not depend on which \(\gamma\) is chosen: for any other
\(\gamma'\in R_{\geq\mathbf{r}}\) we have
\(\gamma'\varphi(\gamma)=\varphi(\gamma'\gamma)=\gamma\varphi(\gamma')\),
whence \(\varphi(\gamma)/\gamma=\varphi(\gamma')/\gamma'\). It therefore
identifies \(C\) with a subring of \(R[1/\gamma]\), and the multiplication
on \(C\) is the one inherited from that identification. This evaluation is
Lemma 5.2.

By the independence just noted the subring obtained does not depend on
which \(\gamma\) is taken; only the localization the computation is
carried out in does. So the one-argument call above makes the choice
itself, taking the first generator of \(R_{\geq\mathbf{r}}\). A second
argument overrides it: \texttt{steinCoordinateAlgebra(homData, i)} takes
the generator in position \texttt{i}, counted from \texttt{0}. The
argument is an integer rather than a ring element because it is that
position and not an arbitrary element of \(R\).

The subring \(A=R_{(0,\geq0)}\) the strand is finite over is not an
argument. A bidegree \((0,n)\) piece of \(S\) is spanned by the monomials
in the target block alone, so \(A\) is generated by the images of the
target-block variables --- the coordinates of \(X\) --- and the degrees of
\texttt{S} already say which variables those are. The point of the whole
computation is that the strand \(C\) is larger than \(A\): the extra
algebra generators are the evaluated Hom generators, and they are what
distinguishes \(Z\) from \(X\).

The returned \texttt{algebraData} holds the graded coordinate algebra of
\(Z\) under \texttt{"ring"}. It is a quotient of a polynomial ring with
one variable for each algebra generator that was used, in this order:
first the target-block variables, then one variable for each evaluated Hom
generator. (The strand generator of bidegree \((0,0)\) is the unit of
\(C\) and gets no variable of its own.) That order is what the printed
output has to be read against: the relation \(p_0p_1-p_2^2\) of Section
\ref{sec:quadratic} has \(p_0,p_1\) the two coordinates of \(X\) and
\(p_2\) the single evaluated Hom generator. The relations themselves are
under \texttt{"definingIdeal"}.

The same \(C\) is also returned in a second form, as a module over \(A\)
rather than as a ring, with its presentation. This is the form in which
\(C\) is finitely generated: a finite presentation over \(A\) is what
makes \(C\) computable at all, and \texttt{pushForward} is what produces
it.

Finally, \texttt{algebraData} keeps \(A\), the localization
\(R[1/\gamma]\), and the \(\gamma\) that defines it. The next step needs
the localization, since that is where the evaluated generators live.

\subsubsection*{The graph ---
  \texorpdfstring{{\normalfont\ttfamily directSteinGraph}}{directSteinGraph}}

The algebra generators of \(Z\) produced by the previous step are
elements \(c_0,\dots,c_N\) of \(C\), homogeneous of bidegrees
\((0,b_0),\dots,(0,b_N)\). 
% They are not functions on \(Y\): a projective
% variety has only the constants for global functions. By the display above
% they are sections of line bundles,
% \[c_i\in C_{b_i}=H^0(Y,f^*\sO_X(b_i))=H^0(Z,g^*\sO_X(b_i)),\]
% realized concretely inside \(R[1/\gamma]\), which is where the
% computation runs.

Each \(c_i\) becomes a homogeneous coordinate \(z_i\) on \(Z\), carrying
the bidegree of the \(c_i\) it stands for:
\[J=S[z_0,\dots,z_N],\qquad \deg z_i=(0,b_i).\]
The map to take the kernel of is
\[\Phi\colon J\longrightarrow R[1/\gamma][t],\qquad
\Phi(x)=\bar x\quad(x\text{ a variable of }S),\qquad
\Phi(z_i)=c_i\,t^{b_i}.\]
Each variable \(x\) of \(S\) goes to its class \(\bar x\) in
\(R=S/I_{\Gamma_f}\), taken inside \(R[1/\gamma]\); that is all the first
half of the map does. The graph of \(h\colon Y\to Z\) is \(\ker\Phi\), a
bihomogeneous ideal of \(J\), and \(J\) is the coordinate ring of the
product of the two spaces containing \(\Gamma_f\) with the weighted
projective space in which the \(z_i\) place \(Z\). It is weighted exactly
when the \(b_i\) are not all equal.

The parameter \(t\) is what forces the answer to be a projective one.
Writing \(F=\sum_\alpha F_\alpha z^\alpha\) with \(F_\alpha\in S\),
\[\Phi(F)=\sum_\alpha
\bar F_\alpha\,c^\alpha\,t^{\langle b,\alpha\rangle},\]
and since \(t\) is a free variable this vanishes if and only if
\(\sum \bar F_\alpha c^\alpha=0\) separately for each value of
\(\langle b,\alpha\rangle=\sum_i b_i\alpha_i\). So \(t\) admits only those
relations that are homogeneous in the \(z\)-degree, which are the ones
that survive the scaling
\(\lambda\cdot(z_0,\dots,z_N)=(\lambda^{b_0}z_0,\dots,\lambda^{b_N}z_N)\)
under which the coordinates of \(Z\) are defined. Without it the kernel
would also record that each \(z_i\) equals one particular element
\(c_i\), which is not a statement about \(Z\) as a projective variety.

Working over \(R[1/\gamma]\) rather than \(R\) is what saturates the
result. The returned \texttt{graphData} holds the kernel under
\texttt{"graphIdeal"}.

\medskip

What the construction requires of \texttt{S} is only that the variables
split into a source block of degrees \((\ast,0)\) and a target block of
degrees \((0,\ast)\): that split is what lets the first degree see the
source direction alone, so that the \((0,\geq0)\)-strand means what it
does. The degrees within a block need not be \(1\), and weighted
projective spaces are allowed, as in the paper. The examples of Section
\ref{sec:examples} are all standard bigraded, but
\texttt{tests/weighted.m2} runs the same construction over \(\PP(1,2)\).

For examples where the minimal resolution is too expensive, the
experimental entry point \texttt{steinHomDataAtBound} accepts a supplied
bound in place of a computed one, and returns the same kind of table with
\texttt{"certifiedBound"} set to false.

\section{Environment and reproducibility}\label{sec:reproducibility}

The computations were run with Macaulay2 1.24.11 on macOS (Apple
Silicon, 2023); that is the version the timings below were measured on
and the version the transcripts above were taken from. The package is
now developed and tested against Macaulay2 1.26.06 on the same machine,
where the examples give the same answers and the timings have not been
measured again. 1.24.11 remains the oldest version the package is
believed to work on. The relevant bundled packages include
\texttt{Truncations}, \texttt{MinimalPrimes}, and \texttt{Saturation}.
Measured wall-clock times on that machine: \texttt{tests/basic.m2} about
4 s, \texttt{tests/global-hom.m2} under 1 s,
\texttt{tests/weighted.m2} under 1 s,
\texttt{tests/mori-fiber-space.m2} about 1 s,
\texttt{tests/blowup-line.m2} about 12 s, so the standard suite takes
roughly 15--20 s in total. The separate twisted-cubic benchmark takes
about 11 s with the supplied bound.

From the project root, the standard suite is run by:

\begin{verbatim}
./run-tests.sh
\end{verbatim}

The slower twisted-cubic benchmark is run separately:

\begin{verbatim}
M2 --no-readline --stop -q tests/blowup-twisted-cubic.m2
\end{verbatim}

The complete input files are included in the repository. The excerpts
below are chosen to show the mathematical input and the checks, without
reproducing every elimination polynomial or the full Macaulay2 startup
log.

\section{Representative Macaulay2 calculations}\label{sec:examples}

\subsection{\texorpdfstring{A quadratic map of \(\PP^1\)}%
                           {A quadratic map of P1}}\label{sec:quadratic}

The first test is the finite map \([s:t]\mapsto[s^2:t^2]\). Its graph is
described by \(y_0^2x_1-y_1^2x_0\). The following excerpt shows the
input and the main output.

\begin{verbatim}
i = ideal(y0^2*x1-y1^2*x0);
d = steinHomData(S,i);
c = steinCoordinateAlgebra(d,0);

d#"bound"
o = {1, 0}

d#"steinGeneratorDegrees"
o = {{0, 0}, {0, 1}}

c#"definingIdeal"
o = ideal(p_0*p_1-p_2^2)

isPrime (directSteinGraph(d,c)#"graphIdeal")
o = true
\end{verbatim}

The ring of the last ideal is generated automatically, so its variables
print as \(p_0,p_1,p_2\); they are the images of \(x_0\), \(x_1\), and
of the evaluated degree-\((0,1)\) Hom generator. Writing them as
\(X_0,X_1,z\) --- which is what \texttt{tests/basic.m2} does, by
rebuilding the map into an explicitly named ring --- the relation is
\(X_0X_1-z^2\).

Thus the Stein intermediate is the source \(\PP^1\), embedded as the conic
\(X_0X_1-z^2=0\): since \(f\) is finite, \(Z\cong Y=\PP^1\), polarized by
\(f^*\sO(1)=\sO(2)\), whose section ring is the second Veronese subring of
\(k[s,t]\).

\subsection{\texorpdfstring{A weighted source \(\PP(1,2)\)}%
                           {A weighted source P(1,2)}}\label{sec:weighted}

Nothing in the construction asks for the degrees inside a block to be
\(1\), so the same code runs over a weighted projective space. Take
\(\PP(1,2)\), with coordinates \(y_0,y_1\) of degrees \((1,0)\) and
\((2,0)\), and the map \([y_0:y_1]\mapsto[y_0^4:y_1^2]\) to \(\PP^1\).
Both coordinates have weighted degree \(4\), so this is a morphism, and
it is finite of degree two.

\begin{verbatim}
S = QQ[y0,y1,x0,x1,Degrees=>{{1,0},{2,0},{0,1},{0,1}}];
i = ideal(x0*y1^2-x1*y0^4);
d = steinHomData(S,i);
c = steinCoordinateAlgebra(d,0);

blockDegreeData S
o = {1, 1, 3, 2}

d#"bound", d#"certifiedBound"
o = ({2, 0}, true)

c#"definingIdeal"
o = ideal(p_0*p_1-p_2^2)

hilbertFunction(1,c#"ring"), hilbertFunction(2,c#"ring")
o = (3, 5)
\end{verbatim}

The third entry of \texttt{blockDegreeData} is what separates this from
the standard case: \(c_1=1+2=3\) is the weight sum, where the variable
count would give \(2\). The distinction is invisible in the standard
examples, and getting it wrong would not fail loudly either --- Corollary
4.3 subtracts \(c_1\), so too small a value only enlarges the bound, and
the answer would stay correct while the computation did more work. That is
why the number is read off the ring instead of being asked for.

Since \(y_0^2\) and \(y_1\) span the weighted degree-two part of
\(k[y_0,y_1]\), they identify \(\PP(1,2)\) with \(\PP^1\), and under that
identification the map is the squaring map. So \(Z\cong\PP^1\) with
\(g^*\sO(1)=\sO(2)\), and the computed ring is the same conic as in
Section \ref{sec:quadratic}, the second Veronese subring of a polynomial
ring in two variables, with \(H(n)=2n+1\).

\subsection{A Mori fiber space with a finite part}\label{sec:mori}

The next example is the composition
\[\PP^1 \times \PP^2 \to \PP^2 \to \PP^2,\]
where the second map is \([x_0:x_1:x_2] \mapsto [x_0^2:x_1^2:x_2^2]\).
The projection has connected \(\PP^1\)-fibers, while the second map is
finite. The relevant part of the test input and output is:

\begin{verbatim}
igraph = kernel graphParametrization;
d = steinHomData(sgraph,igraph);
c = steinCoordinateAlgebra(d,0);

numgens igraph, codim igraph
o = (12, 4)

d#"bound"
o = {2, 0}

dim(c#"ring")
o = 3

hilbertFunction(1,c#"ring"), hilbertFunction(2,c#"ring")
o = (6, 15)

isPrime(c#"definingIdeal")
o = true

fiberComponents = minimalPrimes ifiber;
#fiberComponents
o = 4
\end{verbatim}

Each of the four fiber components is a copy of \(\PP^1\). The computed
Stein intermediate is \(\PP^2\), presented by its second Veronese
subring, and the connected part is the projection to \(\PP^2\).

Here \(Z=\PP^2\) carries the polarization \(g^*\sO(1)=\sO(2)\), so its section
ring is the second Veronese subring of \(k[x_0,x_1,x_2]\), that is, the
subring generated by the degree-2 monomials. Its Hilbert
function is \(H(1)=6\) (the 6 monomials of degree 2 in 3 variables) and
\(H(2)=15\) (the 15 monomials of degree 4 in 3 variables). The
computation reproduces these known invariants, which is what one expects
if the algorithm identifies the Stein intermediate correctly.

\subsection{\texorpdfstring{The blow-up of a line in \(\PP^3\)}%
                           {The blow-up of a line in P3}}\label{sec:blowup-line}

We next take the blow-down \(\Bl_L(\PP^3)\to\PP^3\) along a line
\(L\), followed by the coordinatewise square map on \(\PP^3\). The main
numerical checks are:

\begin{verbatim}
numgens igraph, codim igraph
o = (26, 7)

d = steinHomData(sgraph,igraph);
d#"bound", #d#"steinGeneratorIndices"
o = ({2, 0}, 8)

dim(c#"ring")
o = 4

hilbertFunction(1,c#"ring"), hilbertFunction(2,c#"ring")
o = (10, 35)

isPrime c#"definingIdeal"
o = true

degree(ssource/ie)
o = 2

#fiberPrimes, all(fiberPrimes,p -> degree(rfiber/p) == 1)
o = (8, true)
\end{verbatim}

The Stein ring has the Hilbert values of the second Veronese subring of
\(k[x_0,\dots,x_3]\), as expected for \(Z=\PP^3\) with \(g^*\sO(1)=\sO(2)\): \(H(1)=10\)
(the 10 monomials of degree 2 in 4 variables) and \(H(2)=35\) (the 35
monomials of degree 4 in 4 variables). The computation recovers both
values. Note that the presentation produced by the algorithm uses 11
variables, one of the degree-two module generators being redundant as an
algebra generator; the quotient ring is nevertheless the expected one.

The exceptional divisor of \(\Bl_L(\PP^3) \to \PP^3\) over a line
\(L\) is the projective bundle
\[\PP(N_{L/\PP^3}) = \PP(\sO(1)^{\oplus 2}) \cong \PP^1 \times \PP^1.\]
The computation returns a surface of degree 2, in agreement with the
Segre quadric; the check is carried out on the source side,
independently of the Hom construction.

\subsection{The twisted-cubic blow-up}\label{sec:twisted-cubic}

The final example is the blow-up of \(\PP^3\) along the twisted cubic,
followed by the coordinatewise square map. It is deliberately retained
as a heavier benchmark. Here the source sits in \(\PP^{11}\), so the
automatic Corollary 4.3 bound would require a minimal free resolution
out to homological degree \(11+3+1\); that computation did not finish
within about ten minutes on the machine used here and was aborted. The
test therefore supplies \(\mathbf{r}=(2,0)\) and records that this bound is
experimental; with the supplied bound the example runs in about 11
seconds.

\begin{verbatim}
reesI = kernel map(targetRees,trees,{...});
numgens reesI, codim reesI
o = (3, 2)

igraph = saturate(igraphRaw,ideal matrix{flatten rows});
numgens igraph, codim igraph, isPrime igraph
o = (70, 11, true)

d = steinHomDataAtBound(sgraph,igraph,{2,0});
d#"steinGeneratorDegrees"
o = {{0,0},{0,1},{0,1},{0,1},{0,1},{0,1},{0,1},{0,2}}

hilbertFunction(1,c#"ring"), hilbertFunction(2,c#"ring")
o = (10, 35)

d#"certifiedBound"
o = false
\end{verbatim}

The resulting ring again has the Hilbert values of the second Veronese
subring of \(k[x_0,\dots,x_3]\). One degree-two module generator is algebraically redundant,
so automatic graph construction is not yet completed for this example.
This is a useful illustration of the difference between a bound that
happens to work and a bound certified by Corollary 4.3.

\section{Summary of results}\label{sec:summary}

In the table below, \(\dim\) is the Krull dimension of the computed
Stein coordinate ring, i.e. \(\dim Z+1\) (the ring is the section ring,
so it is the affine cone over \(Z\)).

\begin{longtable}[]{@{}lllll@{}}
\toprule\noalign{}
Example & Bound & \(\dim\) & Selected Hilbert data & Graph \\
\midrule\noalign{}
\endhead
\bottomrule\noalign{}
\endlastfoot
\(\PP^1\) quadratic map & \((1,0)\) & 2 & \(H(1)=3\) & prime \\
\(\PP^1\) cubic map & \((2,0)\) & 2 & twisted cubic & prime \\
\(\PP(1,2)\) degree two & \((2,0)\) & 2 & \(H(1)=3,H(2)=5\) & prime \\
\(\PP^1\times\PP^2\) & \((2,0)\) & 3 & \(H(1)=6,H(2)=15\) & prime \\
\(\Bl_L(\PP^3)\) & \((2,0)\) & 4 & \(H(1)=10,H(2)=35\) & prime \\
Twisted-cubic blow-up & \((2,0)^*\) & 4 & \(H(1)=10,H(2)=35\) &
partial \\
\end{longtable}

The star (*) indicates a supplied experimental bound rather than a bound
computed from the full minimal resolution.

\textbf{Interpretation:} Across these examples, the implementation
reproduces the expected intermediate objects for finite maps,
positive-dimensional connected fibers, a divisorial contraction, and a
weighted projective source. The
graph ideals produced by the direct localization strategy are prime in
the completed examples. The computed Hilbert functions, exceptional
divisors, and fiber structures agree with the classical invariants of
the corresponding geometry. This is evidence that the implementation
behaves as intended on these inputs; it is neither a proof of the
algorithm nor a statement about inputs of a different shape.

\section{Limitations}\label{sec:limitations}

This is a research prototype rather than a general-purpose Macaulay2
package. The main limitations are:

\begin{itemize}
\item
  Deep minimal free resolutions can dominate the exact-bound computation
\item
  Finite module generators may be redundant as algebra generators
\item
  Heuristic stabilization gives evidence but does not prove the explicit
  bound
\item
  The component-selection part of the full algorithm is not automated
  for every possible input
\item
  The weighted examples that have been run are small --- an isomorphism
  and a degree-two map out of \(\PP(1,2)\). Nothing on the scale of the
  blow-up examples has been tried with a weighted grading
\end{itemize}

The calculations should therefore be read as reproducible computational
evidence, not as a replacement for the mathematical hypotheses or proofs
in the source paper.

\section{Use of AI in the implementation}\label{sec:ai}

The Macaulay2 code, the test files, and this note were written
essentially by an AI system (Claude); the author only made minor edits.
The author has read the code and the text and believes them to be
correct, but has not verified every line, every intermediate
computation, or every claim in detail. Readers should keep this in mind.

Two points nevertheless make large errors unlikely. First, the algorithm
being implemented is the one described in \hyperref[sec:references]{Yasuda
(2026)}; nothing conceptually difficult is being attempted here, so the
room for the code to silently do something else is limited. Second, the
tests do not merely report what the code computes: each example is
checked against independently known geometry (Hilbert functions of
Veronese subrings, the degree of the exceptional divisor, the number and
degree of the fiber components, primality of the resulting ideals), and
these checks are run automatically by \texttt{run-tests.sh}. A serious
error in the construction would have to reproduce all of these known
invariants by accident.

None of this replaces a proof, and computations marked as non-certified
are not presented as proofs.

\section{Conclusion}\label{sec:conclusion}

This project shows that the bigraded Hom-module construction can be used
in a practical Macaulay2 workflow on several nontrivial examples. The
most useful outcome is not a large collection of examples, but a
transparent record linking the mathematical construction, the source
code, the actual computer input, and the numerical output. The
repository is therefore intended as a public technical note and
reproducible experiment rather than as a claim of a fully general
implementation.

\section{References}\label{sec:references}

\begin{itemize}
\item
  \textbf{Smith, G. G.} (2000). Computing global extension modules.
  \emph{Journal of Symbolic Computation}, 29(4), 729--746.\\
  (Monograded version of global Hom computation, foundational for the
  bigraded extension.)
\item
  \textbf{Stacks Project Authors.} Stein factorization.\\
  \url{https://stacks.math.columbia.edu/}\\
  (Definition and basic properties of Stein factorization in scheme
  theory.)
\item
  \textbf{Yasuda, T.} (2026). An algorithm for the minimal model program
  in dimension three. arXiv:2603.13703v2.\\
  \url{https://arxiv.org/abs/2603.13703v2}\\
  (Sections 4--5 and Algorithm 1; this repository implements Section
  5.1.)
\item
  \textbf{Grayson, D. R. \& Stillman, M. E.} Macaulay2, a software
  system for research in algebraic geometry.\\
  \url{https://macaulay2.com/}
\end{itemize}

\end{document}
